\documentclass{article}
\usepackage{amsmath}
\begin{document}
\section*{Local gate evidence can precede recovery of a stalled broker}
Keep the same prechecked local optical tube, 10 mm adverse bias,
physically bound gate and actual native resolved courier pose.
Use the same independent diagnostic, timed closure and later
dock receipt. An additional registered team member is neither
sender, receiver nor required relay for these duties. At the
first scheduled clearance dispatch, that member's pose is
unavailable. The pinned native whole-broker pose-update phase
halts with both actual work items still in its FIFO (only the
diagnostic is queued in the local-optical arm). The trace
asserts no qualified callback and no false charge at that cut.
Neither independently named duty is silently discarded.

A separately authorized removal of the unused registration permits
native delivery of the retained items from the original
chronological queue at the later cut. Motion continues normally.
Radio evidence produced at the whole-body threshold does not
become evidence at the gate until the actual callback there.
Removal 30 ms after the first attempt delivers the clearance
about 60.36 / 61.44 ms after outer egress clearance. Removal
60 ms after the attempt yields about 90.36 / 91.44 ms, missing
the 75 ms actuator cutoff at either gate. The diagnosis and later
dock/egress still complete; the upstream radio timestamps are
unchanged. In the matched slow-removal case, a local qualified
optical sample permits the gate effect first (about
40.36 / 41.44 ms); the separate native diagnostic is
retained and actually delivered only after broker recovery.
It does not inherit the gate's actuator deadline.

The observer-to-base diagnostic and its three allowed breadcrumb
sites lie strictly in the outside halfspace of the gate slab, so
none of the candidate hop segments can cross that temporary
slab even when it is installed. This is not a general dynamic
RF propagation model. The independent diagnostic eventually
completes in all arms, the exact late dock duty and audit receipts
complete after those phases, and both successful closure arms
retain the same egress bound. All 24 two-gate / worker-cut /
observer variants agree. Local evidence helps because its causal
availability was not routed through a failed global dependency,
not because the later queue operation is deemed complete early.
This is a held-in example with modeled authorization,
sensing, and interruption timing, not deployment advice or
independent mission acceptance.
\end{document}
